\section*{Chapter \ref{ch:g6:circuits-and-safety} --- Electric Circuits and Safety}

\begin{solution}{exo:g6:circuits-and-safety:1}
Faithfully: the connections --- which device joins which, where
roads fork. Ignored: the real wires' lengths, colors and snaking
paths across the table.
\end{solution}

\begin{solution}{exo:g6:circuits-and-safety:2}
\omterm{def:g3:first-electric-circuit:battery}{Battery}: two unequal parallel bars; lamp: circle with a cross;
open \omterm{ex:g3:first-electric-circuit:switch}{switch}: a lifted drawbridge; closed \omterm{ex:g3:first-electric-circuit:switch}{switch}: the bridge laid
flat; \omterm{ex:g6:circuits-and-safety:motor}{motor}: M in a circle; \omterm{ex:g6:circuits-and-safety:motor}{buzzer}: a struck circle. The
\emph{long} bar marks the \omterm{def:g3:first-electric-circuit:battery}{battery}'s $+$.
\end{solution}

\begin{solution}{exo:g6:circuits-and-safety:3}
One rectangle: \omterm{def:g3:first-electric-circuit:battery}{battery}, \omterm{ex:g3:first-electric-circuit:switch}{push-switch} and \omterm{ex:g6:circuits-and-safety:motor}{buzzer} in a single loop.
Lifting the finger opens the \omterm{ex:g3:first-electric-circuit:switch}{switch} --- a gap in the only road,
and the loop law silences the \omterm{ex:g6:circuits-and-safety:motor}{buzzer} instantly.
\end{solution}

\begin{solution}{exo:g6:circuits-and-safety:4}
Not necessarily. The deciding question: who is connected to whom?
Same devices with the same connections --- same circuit, however
different the tangles look.
\end{solution}

\begin{solution}{exo:g6:circuits-and-safety:5}
The fan spins on: the lamp's death darkens only its own \omterm{def:g5:series-parallel-circuits:parallel}{branch} ---
parallel independence. The open \omterm{ex:g3:first-electric-circuit:switch}{switch} sits on the common road
before the fork, so it stops both: the command-post rule.
\end{solution}

\begin{solution}{exo:g6:circuits-and-safety:6}
A bare conducting road joining the source's two sides directly,
skipping every device: the unhindered current rushes and the wires
\omterm{def:g5:heat-and-insulation:heat}{heat}. Victims: the \omterm{def:g3:first-electric-circuit:battery}{battery} (drained and ruined), the wire's
insulation --- and, on the mains, sometimes the whole house.
\end{solution}

\begin{solution}{exo:g6:circuits-and-safety:7}
It cut its \omterm{def:g5:series-parallel-circuits:parallel}{branch} in a blink because the current began to rush ---
a fault, not a whim. Find and fix the fault (unplug the suspect,
inspect the \omterm{def:g5:series-parallel-circuits:parallel}{branch}) before resetting the guard.
\end{solution}

\begin{solution}{exo:g6:circuits-and-safety:8}
Bath: damp skin conducts, and mains push through a wet body can be
deadly. Damaged cords: cracked insulation is a fence with a gap
--- fingers can meet copper. \omterm{def:g3:first-electric-circuit:battery}{Battery} yes, socket never: the
\omterm{def:g3:first-electric-circuit:battery}{battery}'s push is tiny; the socket's is hundreds of times
mightier.
\end{solution}

\begin{solution}{exo:g6:circuits-and-safety:9}
Diagram: \omterm{def:g3:first-electric-circuit:battery}{battery}, closed switch on the common road, then a fork
--- \omterm{ex:g6:circuits-and-safety:motor}{buzzer} \omterm{def:g5:series-parallel-circuits:parallel}{branch} and lamp \omterm{def:g5:series-parallel-circuits:parallel}{branch} --- rejoining to $-$. Predicted:
the lamp shines \emph{and} the \omterm{ex:g6:circuits-and-safety:motor}{buzzer} \omterm{def:g3:sound-around-us:sound}{sounds}, independently. With
an opened extra switch on the \omterm{ex:g6:circuits-and-safety:motor}{buzzer}'s \omterm{def:g5:series-parallel-circuits:parallel}{branch}: silence there, but
the lamp shines on --- \omterm{def:g5:series-parallel-circuits:parallel}{branch} \omterm{ex:g3:first-electric-circuit:switch}{switches} command only their \omterm{def:g5:series-parallel-circuits:parallel}{branch}.
\end{solution}

\begin{solution}{exo:g6:circuits-and-safety:10}
The extra wire is a \omterm{def:g6:circuits-and-safety:short}{short circuit}: a bare direct road from $+$ to
$-$. Current prefers the easy empty road to the hard-working
\omterm{ex:g6:circuits-and-safety:motor}{motor}, so the fan starves while the shortcut wire carries the rush
and \omterm{def:g5:heat-and-insulation:heat}{heats} --- draining the \omterm{def:g3:first-electric-circuit:battery}{battery} and inviting burns. Remove the
``solidity'' wire.
\end{solution}

\begin{solution}{exo:g6:circuits-and-safety:11}
The wall \omterm{ex:g3:first-electric-circuit:switch}{switch} opens the fixture's own loop --- the lamp's road
is cut. The breaker opens the whole \omterm{def:g5:series-parallel-circuits:parallel}{branch} farther upstream ---
even if the fixture's wiring is faulty or the wall \omterm{ex:g3:first-electric-circuit:switch}{switch}
mislabeled, the \omterm{def:g5:series-parallel-circuits:parallel}{branch} is dead. Two gaps \omterm{def:g4:series-circuits:series}{in series}: either alone
stops the current, both together forgive a mistake about either.
\end{solution}

\begin{solution}{exo:g6:circuits-and-safety:12}
\omterm{def:g3:first-electric-circuit:battery}{Battery}, then the master switch on the common road, then four
parallel \omterm{def:g5:series-parallel-circuits:parallel}{branches}: stage lamp 1; stage lamp 2; curtain \omterm{ex:g6:circuits-and-safety:motor}{motor} with
its own series \omterm{ex:g3:first-electric-circuit:switch}{switch}; \omterm{ex:g6:circuits-and-safety:motor}{buzzer} with its own series \omterm{ex:g3:first-electric-circuit:switch}{switch}. Checks:
master commands all (common road); lamps fail independently
(separate \omterm{def:g5:series-parallel-circuits:parallel}{branches}); \omterm{ex:g6:circuits-and-safety:motor}{motor} and \omterm{ex:g6:circuits-and-safety:motor}{buzzer} each obey their own \omterm{ex:g3:first-electric-circuit:switch}{switch}
(series within their \omterm{def:g5:series-parallel-circuits:parallel}{branch}) and ignore each other (parallel
between \omterm{def:g5:series-parallel-circuits:parallel}{branches}).
\end{solution}

\begin{solution}{pb:g6:circuits-and-safety:1}
\textbf{1.} \omterm{def:g4:series-circuits:series}{In series}, one dead or switched-off device would
darken everything, and three devices would share the push, all
dim and feeble. Parallel keeps each device independent and at
full strength.
\textbf{2.} Each \omterm{ex:g3:first-electric-circuit:switch}{switch} commands exactly its own \omterm{def:g5:series-parallel-circuits:parallel}{branch}'s device;
none of them commands the other \omterm{def:g5:series-parallel-circuits:parallel}{branches}.
\textbf{3.} On the common road, between the \omterm{def:g3:first-electric-circuit:battery}{battery} and the first
fork --- upstream of every \omterm{def:g5:series-parallel-circuits:parallel}{branch}.
\textbf{4.} \omterm{def:g3:first-electric-circuit:battery}{Battery} $\to$ master \omterm{ex:g3:first-electric-circuit:switch}{switch} $\to$ \omterm{def:g5:series-parallel-circuits:parallel}{junction} forking
into three \omterm{def:g5:series-parallel-circuits:parallel}{branches} --- (ceiling lamp $+$ its \omterm{ex:g3:first-electric-circuit:switch}{switch}), (bench
lamp $+$ its \omterm{ex:g3:first-electric-circuit:switch}{switch}), (fan \omterm{ex:g6:circuits-and-safety:motor}{motor} $+$ its \omterm{ex:g3:first-electric-circuit:switch}{switch}) --- rejoining
and returning to the \omterm{def:g3:first-electric-circuit:battery}{battery}.
\textbf{5.} They are \omterm{def:g4:series-circuits:series}{in series} on one \omterm{def:g5:series-parallel-circuits:parallel}{branch}: sharing the push
(both feeble --- the dim ceiling lamp was really the bench pair's
symptom), and the fan's \omterm{ex:g3:first-electric-circuit:switch}{switch}, sitting on their shared road,
commands them both.
\textbf{6.} Split them: run one new wire so the bench lamp gets
its own \omterm{def:g5:series-parallel-circuits:parallel}{branch} across the supply --- each then full-strength and
independent.
\textbf{7.} A \omterm{def:g6:circuits-and-safety:short}{short circuit} across the \omterm{def:g3:first-electric-circuit:battery}{terminals}. Check whether
the \omterm{def:g3:first-electric-circuit:battery}{battery} still delivers --- shorts drain and can ruin cells
(a warm \omterm{def:g3:first-electric-circuit:battery}{battery} deserves retirement).
\textbf{8.} \omterm{def:g6:circuits-and-safety:short}{Short circuits} (bare wires meeting each other) and
shocks or burns (fingers meeting bare copper).
\textbf{9.} Damp skin conducts far better than dry; even the
\omterm{def:g3:first-electric-circuit:battery}{battery} build deserves dry-hands discipline --- and in the house
the same touch meets the mains' push, hundreds of times mightier:
possibly deadly. Habits must be trained on the safe circuit.
\textbf{10.} The breaker detected that too much current was rushing in its
\omterm{def:g5:series-parallel-circuits:parallel}{branch} --- a fault, likely rain in a fitting. Wise order: leave
the breaker off, unplug and inspect the \omterm{def:g5:series-parallel-circuits:parallel}{branch}, fix or dry the
fault, only then reset.
\textbf{11.} Its own loop, separate from the shed's network (or
as one more parallel \omterm{def:g5:series-parallel-circuits:parallel}{branch} on the shed \omterm{def:g3:first-electric-circuit:battery}{battery}): \omterm{def:g3:first-electric-circuit:battery}{battery},
\omterm{ex:g3:first-electric-circuit:switch}{push-switch} at the house, long twin wire, \omterm{ex:g6:circuits-and-safety:motor}{buzzer} at the shed,
back again --- a doorbell is a series loop of its own however far
its wires stretch.
\textbf{12.} For example: ``Sockets: our experiments end at the
\omterm{def:g3:first-electric-circuit:battery}{battery} --- the socket is not ours. Water: wet hands touch no
wiring, indoors or out. Bare wires: every joint wrapped; a bare
wire is a fault, not a shortcut.''
\end{solution}
